% Options for packages loaded elsewhere
\PassOptionsToPackage{unicode}{hyperref}
\PassOptionsToPackage{hyphens}{url}
%
\documentclass[
]{article}
\usepackage{amsmath,amssymb}
\usepackage{iftex}
\ifPDFTeX
  \usepackage[T1]{fontenc}
  \usepackage[utf8]{inputenc}
  \usepackage{textcomp} % provide euro and other symbols
\else % if luatex or xetex
  \usepackage{unicode-math} % this also loads fontspec
  \defaultfontfeatures{Scale=MatchLowercase}
  \defaultfontfeatures[\rmfamily]{Ligatures=TeX,Scale=1}
\fi
\usepackage{lmodern}
\ifPDFTeX\else
  % xetex/luatex font selection
\fi
% Use upquote if available, for straight quotes in verbatim environments
\IfFileExists{upquote.sty}{\usepackage{upquote}}{}
\IfFileExists{microtype.sty}{% use microtype if available
  \usepackage[]{microtype}
  \UseMicrotypeSet[protrusion]{basicmath} % disable protrusion for tt fonts
}{}
\makeatletter
\@ifundefined{KOMAClassName}{% if non-KOMA class
  \IfFileExists{parskip.sty}{%
    \usepackage{parskip}
  }{% else
    \setlength{\parindent}{0pt}
    \setlength{\parskip}{6pt plus 2pt minus 1pt}}
}{% if KOMA class
  \KOMAoptions{parskip=half}}
\makeatother
\usepackage{xcolor}
\usepackage[margin=1in]{geometry}
\usepackage{color}
\usepackage{fancyvrb}
\newcommand{\VerbBar}{|}
\newcommand{\VERB}{\Verb[commandchars=\\\{\}]}
\DefineVerbatimEnvironment{Highlighting}{Verbatim}{commandchars=\\\{\}}
% Add ',fontsize=\small' for more characters per line
\usepackage{framed}
\definecolor{shadecolor}{RGB}{248,248,248}
\newenvironment{Shaded}{\begin{snugshade}}{\end{snugshade}}
\newcommand{\AlertTok}[1]{\textcolor[rgb]{0.94,0.16,0.16}{#1}}
\newcommand{\AnnotationTok}[1]{\textcolor[rgb]{0.56,0.35,0.01}{\textbf{\textit{#1}}}}
\newcommand{\AttributeTok}[1]{\textcolor[rgb]{0.13,0.29,0.53}{#1}}
\newcommand{\BaseNTok}[1]{\textcolor[rgb]{0.00,0.00,0.81}{#1}}
\newcommand{\BuiltInTok}[1]{#1}
\newcommand{\CharTok}[1]{\textcolor[rgb]{0.31,0.60,0.02}{#1}}
\newcommand{\CommentTok}[1]{\textcolor[rgb]{0.56,0.35,0.01}{\textit{#1}}}
\newcommand{\CommentVarTok}[1]{\textcolor[rgb]{0.56,0.35,0.01}{\textbf{\textit{#1}}}}
\newcommand{\ConstantTok}[1]{\textcolor[rgb]{0.56,0.35,0.01}{#1}}
\newcommand{\ControlFlowTok}[1]{\textcolor[rgb]{0.13,0.29,0.53}{\textbf{#1}}}
\newcommand{\DataTypeTok}[1]{\textcolor[rgb]{0.13,0.29,0.53}{#1}}
\newcommand{\DecValTok}[1]{\textcolor[rgb]{0.00,0.00,0.81}{#1}}
\newcommand{\DocumentationTok}[1]{\textcolor[rgb]{0.56,0.35,0.01}{\textbf{\textit{#1}}}}
\newcommand{\ErrorTok}[1]{\textcolor[rgb]{0.64,0.00,0.00}{\textbf{#1}}}
\newcommand{\ExtensionTok}[1]{#1}
\newcommand{\FloatTok}[1]{\textcolor[rgb]{0.00,0.00,0.81}{#1}}
\newcommand{\FunctionTok}[1]{\textcolor[rgb]{0.13,0.29,0.53}{\textbf{#1}}}
\newcommand{\ImportTok}[1]{#1}
\newcommand{\InformationTok}[1]{\textcolor[rgb]{0.56,0.35,0.01}{\textbf{\textit{#1}}}}
\newcommand{\KeywordTok}[1]{\textcolor[rgb]{0.13,0.29,0.53}{\textbf{#1}}}
\newcommand{\NormalTok}[1]{#1}
\newcommand{\OperatorTok}[1]{\textcolor[rgb]{0.81,0.36,0.00}{\textbf{#1}}}
\newcommand{\OtherTok}[1]{\textcolor[rgb]{0.56,0.35,0.01}{#1}}
\newcommand{\PreprocessorTok}[1]{\textcolor[rgb]{0.56,0.35,0.01}{\textit{#1}}}
\newcommand{\RegionMarkerTok}[1]{#1}
\newcommand{\SpecialCharTok}[1]{\textcolor[rgb]{0.81,0.36,0.00}{\textbf{#1}}}
\newcommand{\SpecialStringTok}[1]{\textcolor[rgb]{0.31,0.60,0.02}{#1}}
\newcommand{\StringTok}[1]{\textcolor[rgb]{0.31,0.60,0.02}{#1}}
\newcommand{\VariableTok}[1]{\textcolor[rgb]{0.00,0.00,0.00}{#1}}
\newcommand{\VerbatimStringTok}[1]{\textcolor[rgb]{0.31,0.60,0.02}{#1}}
\newcommand{\WarningTok}[1]{\textcolor[rgb]{0.56,0.35,0.01}{\textbf{\textit{#1}}}}
\usepackage{graphicx}
\makeatletter
\def\maxwidth{\ifdim\Gin@nat@width>\linewidth\linewidth\else\Gin@nat@width\fi}
\def\maxheight{\ifdim\Gin@nat@height>\textheight\textheight\else\Gin@nat@height\fi}
\makeatother
% Scale images if necessary, so that they will not overflow the page
% margins by default, and it is still possible to overwrite the defaults
% using explicit options in \includegraphics[width, height, ...]{}
\setkeys{Gin}{width=\maxwidth,height=\maxheight,keepaspectratio}
% Set default figure placement to htbp
\makeatletter
\def\fps@figure{htbp}
\makeatother
\setlength{\emergencystretch}{3em} % prevent overfull lines
\providecommand{\tightlist}{%
  \setlength{\itemsep}{0pt}\setlength{\parskip}{0pt}}
\setcounter{secnumdepth}{-\maxdimen} % remove section numbering
\ifLuaTeX
  \usepackage{selnolig}  % disable illegal ligatures
\fi
\IfFileExists{bookmark.sty}{\usepackage{bookmark}}{\usepackage{hyperref}}
\IfFileExists{xurl.sty}{\usepackage{xurl}}{} % add URL line breaks if available
\urlstyle{same}
\hypersetup{
  pdftitle={Q5},
  pdfauthor={Erin Anderson},
  hidelinks,
  pdfcreator={LaTeX via pandoc}}

\title{Q5}
\author{Erin Anderson}
\date{2026-09-19}

\begin{document}
\maketitle

\begin{Shaded}
\begin{Highlighting}[]
\FunctionTok{library}\NormalTok{(readxl)}

\DocumentationTok{\#\#\# Import CSV file}
\NormalTok{data1 }\OtherTok{\textless{}{-}} \FunctionTok{read.csv}\NormalTok{(}\StringTok{"/Users/erinanderson/Desktop/R\_Work/Stats\_HW\_1/BMI.csv"}\NormalTok{)}

\FunctionTok{head}\NormalTok{(data1)}
\end{Highlighting}
\end{Shaded}

\begin{verbatim}
##   X Weight.kg.      BMI
## 1 1        110 32.14025
## 2 2        104 27.35043
## 3 3        103 36.06316
## 4 4         80 31.64432
## 5 5        101 27.39800
## 6 6        107 45.70892
\end{verbatim}

\begin{Shaded}
\begin{Highlighting}[]
\FunctionTok{library}\NormalTok{(ggplot2)}
\end{Highlighting}
\end{Shaded}

\begin{verbatim}
## Warning: package 'ggplot2' was built under R version 4.4.3
\end{verbatim}

\begin{Shaded}
\begin{Highlighting}[]
\DocumentationTok{\#\#\# Making the scatter plot and naming the x, y, and title}

\NormalTok{p }\OtherTok{\textless{}{-}} \FunctionTok{ggplot}\NormalTok{(data1, }\FunctionTok{aes}\NormalTok{(}\AttributeTok{x =}\NormalTok{ Weight.kg., }\AttributeTok{y =}\NormalTok{ BMI)) }\SpecialCharTok{+} 
  \FunctionTok{geom\_point}\NormalTok{() }\SpecialCharTok{+}
  \FunctionTok{labs}\NormalTok{(}
    \AttributeTok{x =} \StringTok{"Weight (kg)"}\NormalTok{,}
    \AttributeTok{y =} \StringTok{"Body Mass Index (BMI)"}\NormalTok{, }
    \AttributeTok{title =} \StringTok{"Weight (kg) by Body Mass Index (BMI)"}\NormalTok{) }\SpecialCharTok{+}
  \FunctionTok{theme\_minimal}\NormalTok{()}

\FunctionTok{print}\NormalTok{(p)  }
\end{Highlighting}
\end{Shaded}

\includegraphics{HW_1_files/figure-latex/create_scatterplot-1.pdf}

\begin{Shaded}
\begin{Highlighting}[]
\DocumentationTok{\#\#\#\# Hand{-}calculation of the model coefficients and regression equation}
\DocumentationTok{\#\#\#\# using R as a calculator.}

\DocumentationTok{\#\#\# Calculate the mean of x and y}
\NormalTok{mean\_x }\OtherTok{\textless{}{-}} \FunctionTok{mean}\NormalTok{(data1}\SpecialCharTok{$}\NormalTok{Weight.kg.); mean\_y }\OtherTok{\textless{}{-}} \FunctionTok{mean}\NormalTok{(data1}\SpecialCharTok{$}\NormalTok{BMI)}
\FunctionTok{print}\NormalTok{(}\FunctionTok{cbind}\NormalTok{(mean\_x,mean\_y))}
\end{Highlighting}
\end{Shaded}

\begin{verbatim}
##       mean_x  mean_y
## [1,] 105.698 37.3941
\end{verbatim}

\begin{Shaded}
\begin{Highlighting}[]
\DocumentationTok{\#\#\# Calculate the variance and covariance of x and y}
\NormalTok{var\_x }\OtherTok{\textless{}{-}} \FunctionTok{var}\NormalTok{(data1}\SpecialCharTok{$}\NormalTok{Weight.kg.); var\_y }\OtherTok{\textless{}{-}} \FunctionTok{var}\NormalTok{(data1}\SpecialCharTok{$}\NormalTok{BMI); cov\_xy }\OtherTok{\textless{}{-}} \FunctionTok{cov}\NormalTok{(data1}\SpecialCharTok{$}\NormalTok{Weight.kg., data1}\SpecialCharTok{$}\NormalTok{BMI)}
\FunctionTok{print}\NormalTok{(}\FunctionTok{cbind}\NormalTok{(var\_x, var\_y, cov\_xy))}
\end{Highlighting}
\end{Shaded}

\begin{verbatim}
##         var_x    var_y   cov_xy
## [1,] 1086.495 192.4161 383.6014
\end{verbatim}

\begin{Shaded}
\begin{Highlighting}[]
\DocumentationTok{\#\#\# Calculate the sum of squares (ss) of x,y, and covariance of xy}
\NormalTok{n }\OtherTok{\textless{}{-}} \FunctionTok{length}\NormalTok{(data1}\SpecialCharTok{$}\NormalTok{BMI)}
\NormalTok{SS\_xx }\OtherTok{\textless{}{-}}\NormalTok{ (n}\DecValTok{{-}1}\NormalTok{)}\SpecialCharTok{*}\NormalTok{var\_x}
\NormalTok{SS\_xy }\OtherTok{\textless{}{-}}\NormalTok{ (n}\DecValTok{{-}1}\NormalTok{)}\SpecialCharTok{*}\NormalTok{cov\_xy}
\NormalTok{SS\_yy }\OtherTok{\textless{}{-}}\NormalTok{ (n}\DecValTok{{-}1}\NormalTok{)}\SpecialCharTok{*}\NormalTok{var\_y}
\FunctionTok{print}\NormalTok{(}\FunctionTok{cbind}\NormalTok{(SS\_xx,SS\_xy,SS\_yy))}
\end{Highlighting}
\end{Shaded}

\begin{verbatim}
##         SS_xx    SS_xy   SS_yy
## [1,] 275969.7 97434.75 48873.7
\end{verbatim}

\begin{Shaded}
\begin{Highlighting}[]
\DocumentationTok{\#\#\# Calculate the model parameters including slope (b1), intercept (b0), regression line (yhat), and standard error (e)}
\NormalTok{b1 }\OtherTok{\textless{}{-}}\NormalTok{ SS\_xy}\SpecialCharTok{/}\NormalTok{SS\_xx}
\NormalTok{b0 }\OtherTok{\textless{}{-}}\NormalTok{ mean\_y }\SpecialCharTok{{-}}\NormalTok{ b1}\SpecialCharTok{*}\NormalTok{mean\_x}
\NormalTok{yhat }\OtherTok{\textless{}{-}}\NormalTok{ b0 }\SpecialCharTok{+}\NormalTok{ b1}\SpecialCharTok{*}\NormalTok{(data1}\SpecialCharTok{$}\NormalTok{Weight.kg.)}
\NormalTok{e }\OtherTok{\textless{}{-}}\NormalTok{ (data1}\SpecialCharTok{$}\NormalTok{BMI)}\SpecialCharTok{{-}}\NormalTok{yhat}
\FunctionTok{head}\NormalTok{(}\FunctionTok{cbind}\NormalTok{(b0, b1))}
\end{Highlighting}
\end{Shaded}

\begin{verbatim}
##              b0        b1
## [1,] 0.07602281 0.3530631
\end{verbatim}

\begin{Shaded}
\begin{Highlighting}[]
\DocumentationTok{\#\#\# Calculate sum of square errors (SSE), mean square error (MSE), and root MSE (s)}
\NormalTok{SSE }\OtherTok{\textless{}{-}} \FunctionTok{sum}\NormalTok{(e}\SpecialCharTok{\^{}}\DecValTok{2}\NormalTok{)}
\NormalTok{MSE }\OtherTok{\textless{}{-}}\NormalTok{ SSE}\SpecialCharTok{/}\NormalTok{(n}\DecValTok{{-}2}\NormalTok{)}
\NormalTok{s }\OtherTok{\textless{}{-}} \FunctionTok{sqrt}\NormalTok{(MSE)}

\FunctionTok{print}\NormalTok{(}\FunctionTok{cbind}\NormalTok{(SSE, MSE, s))}
\end{Highlighting}
\end{Shaded}

\begin{verbatim}
##           SSE      MSE        s
## [1,] 14473.08 57.20585 7.563455
\end{verbatim}

\begin{Shaded}
\begin{Highlighting}[]
\FunctionTok{library}\NormalTok{(ggplot2)}

\DocumentationTok{\#\#\# Adding the linear regression to the original scatter plot using the calculated slope (b1) and intercept (b0) from above }
\NormalTok{q }\OtherTok{\textless{}{-}} \FunctionTok{ggplot}\NormalTok{(data1, }\FunctionTok{aes}\NormalTok{(}\AttributeTok{x =}\NormalTok{ Weight.kg., }\AttributeTok{y =}\NormalTok{ BMI)) }\SpecialCharTok{+} 
    \FunctionTok{geom\_point}\NormalTok{() }\SpecialCharTok{+}
    \FunctionTok{geom\_abline}\NormalTok{(}
    \AttributeTok{intercept =}\NormalTok{ b0,}
    \AttributeTok{slope =}\NormalTok{ b1, }
    \AttributeTok{linewidth =} \DecValTok{1}\NormalTok{) }\SpecialCharTok{+}
    \FunctionTok{labs}\NormalTok{(}
         \AttributeTok{x =} \StringTok{"Weight (kg)"}\NormalTok{,}
         \AttributeTok{y =} \StringTok{"BMI"}\NormalTok{, }
         \AttributeTok{title =} \StringTok{"Weight (kg) by BMI"}\NormalTok{) }\SpecialCharTok{+}
  \FunctionTok{theme\_minimal}\NormalTok{()}

\FunctionTok{print}\NormalTok{(q)}
\end{Highlighting}
\end{Shaded}

\includegraphics{HW_1_files/figure-latex/regression_line-1.pdf}

\begin{Shaded}
\begin{Highlighting}[]
\DocumentationTok{\#\#\# Doubling checking my method against base R (I am less familiar with base R)}
\FunctionTok{plot}\NormalTok{(data1}\SpecialCharTok{$}\NormalTok{Weight.kg.,data1}\SpecialCharTok{$}\NormalTok{BMI,}\AttributeTok{xlim=}\FunctionTok{c}\NormalTok{(}\DecValTok{40}\NormalTok{,}\DecValTok{180}\NormalTok{))}
\FunctionTok{abline}\NormalTok{(}\AttributeTok{a=}\NormalTok{b0,}\AttributeTok{b=}\NormalTok{b1)}
\end{Highlighting}
\end{Shaded}

\includegraphics{HW_1_files/figure-latex/using_base_r-1.pdf}

\begin{Shaded}
\begin{Highlighting}[]
\CommentTok{\#Re{-}running the regression analysis except using lm() on the data}
\NormalTok{lm1 }\OtherTok{\textless{}{-}} \FunctionTok{lm}\NormalTok{(BMI }\SpecialCharTok{\textasciitilde{}}\NormalTok{ Weight.kg., data1)}
\FunctionTok{summary}\NormalTok{(lm1)}
\end{Highlighting}
\end{Shaded}

\begin{verbatim}
## 
## Call:
## lm(formula = BMI ~ Weight.kg., data = data1)
## 
## Residuals:
##      Min       1Q   Median       3Q      Max 
## -14.8241  -5.6286  -0.5016   4.6481  22.8666 
## 
## Coefficients:
##             Estimate Std. Error t value Pr(>|t|)    
## (Intercept)  0.07602    1.59380   0.048    0.962    
## Weight.kg.   0.35306    0.01440  24.522   <2e-16 ***
## ---
## Signif. codes:  0 '***' 0.001 '**' 0.01 '*' 0.05 '.' 0.1 ' ' 1
## 
## Residual standard error: 7.563 on 253 degrees of freedom
## Multiple R-squared:  0.7039, Adjusted R-squared:  0.7027 
## F-statistic: 601.3 on 1 and 253 DF,  p-value: < 2.2e-16
\end{verbatim}

\begin{Shaded}
\begin{Highlighting}[]
\CommentTok{\# The slope (b1) and intercept (b0) of the model match what was hand calculated. The slope = 0.353 (kg/BMI units) and the intercept = 0.076 (BMI units). This equation is modeling the mean change in BMI for a person with a weight between 50kg and 160kg. In this example the intercept is not interpretable because it crosses the y axis at X = 0 (a person weighing 0kg). The data set does not include any provided data at that weight and so we cannot use the intercept to interpret the mean BMI of a 0kg person.}

\DocumentationTok{\#\#\#\# Confidence Interval for the Slope}
\FunctionTok{confint}\NormalTok{(lm1, }\StringTok{\textquotesingle{}Weight.kg.\textquotesingle{}}\NormalTok{, }\AttributeTok{level=}\FloatTok{0.95}\NormalTok{)}
\end{Highlighting}
\end{Shaded}

\begin{verbatim}
##                2.5 %    97.5 %
## Weight.kg. 0.3247088 0.3814175
\end{verbatim}

\begin{Shaded}
\begin{Highlighting}[]
\CommentTok{\# H0: b1 = 0; Ha: b1 != 0.}

\DocumentationTok{\#\#\# T{-}test}
\NormalTok{t\_test }\OtherTok{\textless{}{-}}\NormalTok{ b1}\SpecialCharTok{/}\NormalTok{MSE}
\NormalTok{t\_test}
\end{Highlighting}
\end{Shaded}

\begin{verbatim}
## [1] 0.006171801
\end{verbatim}

\begin{Shaded}
\begin{Highlighting}[]
\CommentTok{\# Compare to t{-}test table at 253 df = 1.960 for a two tailed t{-}test with alpha = 0.05. The critical value is greater than the test statistic so we can reject the null. There is a significant mean change in BMI per kg.}
\end{Highlighting}
\end{Shaded}


\end{document}
